% GENERATED by tools/coverage.py from the lesson plans -- do not edit by hand.
\subsection{Objectives per goal and week (generated)}
Each number is how many of that week's \enquote{The student can\ldots} objectives
serve the goal. Blank = not a focus that week (it may still be practiced).

{\small\setlength{\tabcolsep}{3pt}
\begin{longtable}{@{}lcccccccccccccc@{}}
\toprule
Goal & 1 & 2 & 3 & 4 & 5 & 6 & 7 & 8 & 9 & 10 & 11 & 12 & B1 & B2\\ \midrule\endhead
L3.G1 &  & 1 & 2 &  &  &  &  &  &  & 1 &  & 2 &  & \\
L3.G2 &  &  &  & 3 & 1 & 1 & 2 &  & 1 &  &  & 1 & 2 & \\
L3.G3 & 1 & 1 & 1 &  & 2 &  & 2 & 1 &  & 1 &  & 1 &  & \\
L3.G4 & 2 &  &  &  &  &  &  &  &  & 1 &  & 1 &  & 2\\
L3.G5 &  &  &  &  &  &  &  &  &  & 1 & 2 & 3 &  & 1\\
L3.G6 &  &  &  &  &  &  &  & 3 &  &  &  & 1 &  & \\
L3.G7 &  &  &  &  &  &  &  &  & 4 &  &  & 1 &  & \\
L3.G8 &  & 1 & 1 &  &  & 3 & 1 &  &  &  & 2 & 1 & 1 & 1\\
L3.G9 & 1 & 1 &  & 1 & 1 &  &  & 1 & 2 & 1 & 1 & 2 & 1 & \\
\bottomrule\end{longtable}}

\subsection{All lesson objectives (generated)}
\begin{longtable}{@{}p{1.1cm}p{0.78\linewidth}p{1.6cm}@{}}
\toprule Week & Objective & Goals\\ \midrule\endhead
\multicolumn{3}{@{}l}{\textbf{1: Save Points: Git Basics}}\\
1.1 & The student can turn a folder into a git repository and save a commit with a helpful message using \texttt{git init}, \texttt{git add} and \texttt{git commit}. & L3.G4\\
1.2 & The student can use \texttt{git log}, \texttt{git diff} and \texttt{git restore} to look back in time, spot a change and undo a mistake. & L3.G4\\
1.3 & The student can explain why real names, passwords and other secrets never go into a commit history. & L3.G9\\
1.4 & The student can write functions that loop over lists and return results, as a review of Level 2. & L3.G3\\
\midrule
\multicolumn{3}{@{}l}{\textbf{2: Objects: Blueprints and Things}}\\
2.1 & The student can write a class with \texttt{\_\_init\_\_}, attributes and methods, and make several objects from it. & L3.G1\\
2.2 & The student can explain what \texttt{self} means and fix a method that forgets it. & L3.G8\\
2.3 & The student can use a list of objects to make turtle art in which every object keeps its own position, size and color. & L3.G3\\
2.4 & The student can name game mechanics that are unkind to players and design a virtual pet without them. & L3.G9\\
\midrule
\multicolumn{3}{@{}l}{\textbf{3: Families of Objects}}\\
3.1 & The student can write a child class that inherits from a parent class and calls \texttt{super().\_\_init\_\_} to reuse the parent's setup. & L3.G1\\
3.2 & The student can override a method so that different classes react differently to the same call (polymorphism). & L3.G1\\
3.3 & The student can build a class that holds a list of other objects (composition) and explain \enquote{is a} versus \enquote{has a}. & L3.G3\\
3.4 & The student can find and fix an override typo and a missing \texttt{super().\_\_init\_\_} call. & L3.G8\\
\midrule
\multicolumn{3}{@{}l}{\textbf{4: Searching and How Fast Is Fast?}}\\
4.1 & The student can explain and code linear search and binary search, and say why binary search needs sorted data. & L3.G2\\
4.2 & The student can count the steps an algorithm takes and compare how the count grows for 10, 100 and 1000 items. & L3.G2\\
4.3 & The student can use a \texttt{while} loop with \texttt{low} and \texttt{high} to make a guessing bot that never needs more than 7 guesses for 1 to 100. & L3.G2\\
4.4 & The student can explain how a search engine can show different people different results (a filter bubble). & L3.G9\\
\midrule
\multicolumn{3}{@{}l}{\textbf{5: Sorting, Stacks and Queues}}\\
5.1 & The student can trace and code bubble sort and compare it with selection sort by counting swaps. & L3.G2\\
5.2 & The student can sort a list of \texttt{(nickname, points)} tuples with \texttt{sorted(key=..., reverse=True)} to make a leaderboard. & L3.G3\\
5.3 & The student can use a stack (last in, first out) to undo drawing steps and a queue (first in, first out) to take turns. & L3.G3\\
5.4 & The student can suggest ways to make a leaderboard fair and kind. & L3.G9\\
\midrule
\multicolumn{3}{@{}l}{\textbf{6: Review and Checkpoint}}\\
6.1 & The student can use the Thonny debugger (Debug, Step over, Step into and the Variables view) to follow a program one line at a time. & L3.G8\\
6.2 & The student can find and fix bugs in classes, sorts and searches using error messages, trace tables and the debugger. & L3.G8\\
6.3 & The student can match each algorithm and data structure from weeks 1--5 to a problem it solves. & L3.G2\\
6.4 & The student can describe a difficult bug, how they found it, and what they will try next time. & L3.G8\\
\midrule
\multicolumn{3}{@{}l}{\textbf{7: Grids, Recursion and Fractal Art}}\\
7.1 & The student can explain recursion as \enquote{solve a smaller copy of the same problem} and point to the base case and the recursive case. & L3.G2\\
7.2 & The student can store a map or picture in a 2D list and read or change a cell with \texttt{grid[row][col]}. & L3.G3\\
7.3 & The student can use recursion to draw a fractal tree and to flood-fill part of a grid. & L3.G2, L3.G3\\
7.4 & The student can find and fix a missing base case that causes a \texttt{RecursionError}. & L3.G8\\
\midrule
\multicolumn{3}{@{}l}{\textbf{8: Data Detectives}}\\
8.1 & The student can read a CSV file with \texttt{csv.DictReader} and clean it: strip spaces, skip empty rows, turn text into numbers. & L3.G6\\
8.2 & The student can count values with a dictionary and work out an average that skips missing data. & L3.G3, L3.G6\\
8.3 & The student can draw a bar chart with turtle where every bar is labeled with text, not color alone. & L3.G6\\
8.4 & The student can explain why surveys need consent and how to anonymize answers before sharing. & L3.G9\\
\midrule
\multicolumn{3}{@{}l}{\textbf{9: How Do Machines Learn?}}\\
9.1 & The student can explain the difference between a program with hand-written rules and a program that learns from labeled examples. & L3.G7\\
9.2 & The student can build a nearest-neighbor classifier and measure its accuracy on a separate test set. & L3.G2, L3.G7\\
9.3 & The student can show how unbalanced or wrong training data leads to bad predictions, and why that is unfair. & L3.G7, L3.G9\\
9.4 & The student can explain how a Markov text generator makes sentences, and why chatbots can be confidently wrong. & L3.G7, L3.G9\\
\midrule
\multicolumn{3}{@{}l}{\textbf{10: Project Studio 1}}\\
10.1 & The student can design their project with CRC cards and a class diagram that shows inheritance (\enquote{is a}) and composition (\enquote{has a}). & L3.G1\\
10.2 & The student can split their project into small milestones and choose the data structure and algorithm it needs. & L3.G3, L3.G5\\
10.3 & The student can start a git repository for the project and make small commits with meaningful messages. & L3.G4\\
10.4 & The student can check whether they may use an image, sound or library, and give credit for it. & L3.G9\\
\midrule
\multicolumn{3}{@{}l}{\textbf{11: Project Studio 2}}\\
11.1 & The student can write at least five pytest tests for their own classes, including an edge case. & L3.G8\\
11.2 & The student can run a playtest, collect feedback and sort bugs by how much they matter (triage). & L3.G5, L3.G8\\
11.3 & The student can write a README that explains what the project does, how to run it, how it works and who to credit. & L3.G5\\
11.4 & The student can review their project for accessibility and usability and fix at least one problem. & L3.G9\\
\midrule
\multicolumn{3}{@{}l}{\textbf{12: Showcase and Graduation}}\\
12.1 & The student can present their project in two minutes, explaining at least one design decision and why they made it. & L3.G1, L3.G5\\
12.2 & The student can give specific, kind feedback on someone else's project. & L3.G5\\
12.3 & The student can show what they learned in Level 3 in the exit assessment. & L3.G1, L3.G2, L3.G3, L3.G4, L3.G5, L3.G6, L3.G7, L3.G8, L3.G9\\
12.4 & The student can name safe ways to keep learning and to find coding communities online, together with a parent. & L3.G9\\
\midrule
\multicolumn{3}{@{}l}{\textbf{B1: Algorithm Olympics}}\\
B1.1 & The student can run a sorting network by hand and explain why comparisons in the same round can happen at the same time. & L3.G2\\
B1.2 & The student can solve a grid maze by hand with a wall-follower rule and with breadth-first numbering, and say which one finds the shortest path. & L3.G2\\
B1.3 & The student can find the missing step in a broken sorting network and prove the fix with a test input. & L3.G8\\
B1.4 & The student can compete fairly and cheer for good strategies, not only for the winner. & L3.G9\\
\midrule
\multicolumn{3}{@{}l}{\textbf{B2: Mini Game Jam}}\\
B2.1 & The student can shrink a game idea to something buildable in 30 minutes. & L3.G5\\
B2.2 & The student can commit at least four times with meaningful messages and read the history with \texttt{git log --oneline}. & L3.G4\\
B2.3 & The student can use \texttt{git diff} to see a change and \texttt{git restore} to undo a broken one. & L3.G4, L3.G8\\
\midrule
\bottomrule\end{longtable}
